\documentclass[11pt,letterpaper]{article}
\usepackage[margin=0.82in,headheight=14pt,headsep=15pt]{geometry}
\usepackage[T1]{fontenc}
\usepackage{lmodern,amsmath,amssymb,amsthm,mathtools,microtype}
\usepackage{booktabs,tabularx,array,enumitem,xcolor,fancyhdr}
\definecolor{inkblue}{RGB}{24,58,88}
\usepackage[colorlinks=true,allcolors=inkblue]{hyperref}
\hypersetup{pdfpagelayout=SinglePage,pdfpagemode=UseNone,pdfstartview=Fit,
pdftitle={Springer--Borel images in classical type: from Bruhat geometry to equivariant K-theory},
pdfsubject={Published foundations, integral image deductions, and a classical equivariant research program}}
\pagestyle{fancy}\fancyhf{}
\fancyhead[L]{\small Springer--Borel images in classical type}
\fancyhead[R]{\small Expository research note}
\fancyfoot[L]{\footnotesize 22 October 2026}\fancyfoot[R]{\thepage}
\renewcommand{\headrulewidth}{0.3pt}
\setlength{\parindent}{1em}\setlength{\parskip}{3pt}
\setlength{\emergencystretch}{2em}
\setlength{\abovedisplayskip}{7pt plus 2pt minus 2pt}
\setlength{\belowdisplayskip}{7pt plus 2pt minus 2pt}
\setlist{nosep,leftmargin=1.5em}
\newcommand{\Z}{\mathbb Z}\newcommand{\Q}{\mathbb Q}\newcommand{\C}{\mathbb C}
\newcommand{\Sym}{\operatorname{Sym}}\newcommand{\Rep}{\operatorname{Rep}}
\newcommand{\SL}{\operatorname{SL}}\newcommand{\GL}{\operatorname{GL}}
\newcommand{\SO}{\operatorname{SO}}\newcommand{\Sp}{\operatorname{Sp}}
\newcommand{\Spin}{\operatorname{Spin}}\newcommand{\IG}{\operatorname{IG}}
\newcommand{\OG}{\operatorname{OG}}\newcommand{\Gr}{\operatorname{Gr}}
\newcommand{\Sat}{\operatorname{Sat}}\newcommand{\Fitt}{\operatorname{Fitt}}
\newcommand{\ch}{\operatorname{ch}}\newcommand{\im}{\operatorname{im}}
\newcommand{\rk}{\operatorname{rank}}\newcommand{\Tor}{\operatorname{Tor}}
\newcommand{\eps}{\epsilon}\newcommand{\poly}{\mathrm{poly}}
\newcommand{\qb}[3]{\genfrac{[}{]}{0pt}{}{#1}{#2}_{#3}}
\newcolumntype{Y}{>{\raggedright\arraybackslash}X}
\newtheorem{theorem}{Theorem}[section]
\newtheorem{proposition}[theorem]{Proposition}
\newtheorem{example}[theorem]{Example}
\theoremstyle{definition}\newtheorem{definition}[theorem]{Definition}
\newcommand{\record}[2]{\href{https://github.com/sean-rogers-t/Sage-programs/blob/2bc20702f142d6124a22ab97aaa511caf227f916/borel/#1}{#2}}
\begin{document}
\begin{center}
{\LARGE\bfseries Springer--Borel images in classical type}\\[6pt]
{\Large From Bruhat geometry to equivariant K-theory}\\[7pt]
{\small An exposition of the published maps, integral refinements, and a research program}
\end{center}
\begin{abstract}
Springer morphisms turn a representation of a reductive group into a
polynomial subring of its Levi representation rings. The Borel
characteristic map then sends these polynomial characters to flag
cohomology. We explain the geometry and algebra behind this construction,
including Bruhat cells, Schubert bases, Weyl invariants, the meaning of the
classical parameters, and the distinction between normalized and literal
integral sources. We review the published results of Kumar,
Kumar--Rogers, and Kumar--Xie, and the subsequent internally audited
integral-image deductions. We then introduce equivariant cohomology,
ordinary and equivariant K-theory, and the precise image problems these
theories suggest. The immediate emphasis is on uniform results in the
classical families \(A,B,C,D\).
\end{abstract}
\noindent\textbf{Reading guide and status.}
Sections~1--5 develop the background and published construction.
Section~6 explains the integral refinements; Sections~7--9 develop the
new targets and the available comparison results. Section~10 states the
remaining program. A statement labeled a \emph{repository deduction}
means that a proof is written in the cited companion notes and has been
checked internally. It does not mean external peer review or an
established claim of originality. Proposed extensions are stated as
problems, not as completed theorems.

\section{Groups, weights, and the meaning of a parabolic}
Let \(G\) be a connected complex reductive group, \(T\subset B\subset G\)
a maximal torus and a Borel subgroup, and \(P\supset B\) a standard
parabolic. Let \(L\supset T\) be its Levi subgroup. We write
\[
 X=X(T)=\operatorname{Hom}(T,\C^\times),\qquad
 W=N_G(T)/T,\qquad W_L=N_L(T)/T.
\]
A weight \(\lambda\in X\) is a multiplicative character of \(T\).
Its differential is a linear function on
\(\mathfrak t=\operatorname{Lie}(T)\). Thus the same weight gives both
a function \(e^\lambda\) on the torus and a linear function \(\lambda\)
on its Lie algebra. Keeping these two functions distinct is essential.
The roots describe the conjugation action on \(\mathfrak g\); the
simple roots \(\alpha_i\) determine the simple reflections \(s_i\).
The fundamental weights \(\omega_j\) satisfy
\(\langle\omega_j,\alpha_i^\vee\rangle=\delta_{ij}\).

For most integral comparisons we take \(G\) simple and simply connected.
Then \(X\) is the full weight lattice. The root lattice
\(Q_{\rm rt}=\sum_i\Z\alpha_i\) is a finite-index sublattice of \(X\).
The Lie algebra and Weyl group do not determine this character lattice:
\(\Spin_N\) and \(\SO_N\), for example, have different character lattices.
Their complex flag varieties can agree while their specified
representation sources and equivariant coefficient rings differ.

A standard parabolic is specified by its \emph{retained} simple nodes,
which generate \(W_L\). The omitted nodes specify the steps of a
partial flag. One omitted node means a maximal parabolic, several
omitted nodes mean a more general partial flag, all nodes omitted
give \(G/B\), and none omitted give a point. The word parabolic
therefore does not restrict the discussion to Grassmannians.

In type \(A_{m-1}\), take \(G=\SL_m\). The quotient for steps
\(0<r_1<\cdots<r_s<m\) parametrizes chains
\[
 U_{r_1}\subset\cdots\subset U_{r_s}\subset\C^m,\qquad
 \dim U_{r_j}=r_j.
\]
The Levi block sizes are \(r_1,r_2-r_1,\ldots,m-r_s\), and its Weyl
group permutes coordinates within each block. For one step \(r\),
the variety is \(\Gr(r,m)\), of complex dimension \(r(m-r)\).

The notation in Parts II and III instead emphasizes the size of a
remaining classical factor. Put \(a=n-k\). The parameter dictionary is
\[
\begin{array}{c|c|c|c}
\text{type}&\text{published group}&G/P_{n-k}&\text{Levi factor}\\ \hline
C_n&\Sp(2n)&\IG(n-k,2n)&\GL(n-k)\times\Sp(2k)\\
B_n&\SO(2n+1)&\OG(n-k,2n+1)&\GL(n-k)\times\SO(2k+1)\\
D_n&\SO(2n)&\OG(n-k,2n)&\GL(n-k)\times\SO(2k)
\end{array}
\]
Here \(\IG\) uses an alternating form and \(\OG\) a symmetric form;
the subspace must be isotropic for that form. Thus \(n\) is the rank,
\(n-k=a\) is the isotropic subspace dimension, and \(k\) is the residual
classical rank. In the stated published ranges \(1\le k<n\) in B/C
and \(n\ge4,\ 2\le k<n\) in D. Maximal isotropic varieties require
their own boundary conventions; in even orthogonal type there are
two maximal isotropic components.

\section{Bruhat decomposition: where the cohomology basis comes from}
The Bruhat decomposition is the finite stratification
\[
 G/P=\coprod_{w\in W^P} C_w,\qquad
 C_w=BwP/P\simeq\C^{\ell(w)}.
 \tag{1}
\]
Here \(W^P\) is the set of minimal representatives for the right cosets
\(W/W_L\), and \(\ell(w)\) is Coxeter length, the minimum number of
simple reflections in an expression for \(w\). Minimality is equivalent
to having no right descent in a retained Levi node. The closure
relations among the cells define Bruhat order:
\(\overline{C_w}=\bigcup_{v\le w}C_v\).
An algebraic filtration by these closures gives a cellular description
with only even real dimensions.

There is also the opposite decomposition for the opposite Borel
\(B^-\). The opposite Schubert variety
\(X^w=\overline{B^-wP/P}\) has complex codimension \(\ell(w)\).
We use the notation of the Kumar papers:
\[
 \eps_w=[X^w]\in H^{2\ell(w)}(G/P;\Z),\qquad w\in W^P.
 \tag{2}
\]
These classes form an integral basis. In particular, odd cohomology
vanishes and even cohomology is torsion-free. The generating polynomial
\(\sum_{w\in W^P}z^{\ell(w)}\) counts the ranks of \(H^{2d}\).
Torsion in a quotient by a selected image is consequently a property
of that image, not torsion in flag cohomology itself.

For \(G=\SL_3\) and \(P=B\), the six indices have lengths
\[
\begin{array}{c|cccc}
\ell&0&1&2&3\\ \hline
w&1&s_1,s_2&s_1s_2,s_2s_1&s_1s_2s_1
\end{array}
\]
and the rank polynomial is \(1+2z+2z^2+z^3\).
For \(\Gr(r,m)\), the same indices can be encoded by partitions
\(\lambda\subseteq(m-r)^r\): there are at most \(r\) parts, each at
most \(m-r\), and \(\eps_\lambda\) lies in degree \(2|\lambda|\).
Partitions are thus finite-dimensional geometric coordinates for the
Schubert basis, not parameters of the ambient group.

This basis is the bridge to exact computation. A characteristic
polynomial is expanded in the \(\eps_w\), producing a column of integers
or rationals. Products of specified source generators give a matrix
whose column lattice is exactly the image we wish to understand.
The use of Schubert coordinates is geometric, not an arbitrary
choice of a numerical basis. Background on these decompositions and
their compatibility with parabolic projections is given in
Brion's lectures~\cite{Brion}.

\section{Representation rings and two kinds of variables}
The representation ring \(R(L)=\Rep(L)\) is the Grothendieck ring of
finite-dimensional algebraic representations: direct sum gives addition
and tensor product gives multiplication. Its elements are virtual
representations, so negative coefficients are allowed. Restriction of
characters to \(T\) identifies it with
\[
 R(T)=\Z[X],\qquad R(L)=\Z[X]^{W_L},\qquad R(G)=\Z[X]^W.
 \tag{3}
\]
The group ring has basis \(e^\lambda\), with
\(e^\lambda e^\mu=e^{\lambda+\mu}\).
Choosing a lattice basis makes it a Laurent polynomial ring.
For a simply connected semisimple \(G\), \(R(G)\) is a polynomial
ring in the fundamental irreducible characters. This does not say
that every Levi ring is a polynomial ring in arbitrary torus variables.

For example, on \(\GL_r\) let \(t_i=e^{\eps_i}\). Then
\[
 R(\GL_r)=
 \Z[t_1^{\pm1},\ldots,t_r^{\pm1}]^{S_r}
 =\Z[e_1(t),\ldots,e_r(t),e_r(t)^{-1}].
 \tag{4}
\]
The generator \(e_i(t)\) is the character of the \(i\)-th exterior
power, and \(e_r(t)\) is the determinant character. The ordinary
polynomial representation ring omits the determinant inverse.
For \(\SL_m\) the torus relation is \(\prod_i t_i=1\).

Irreducible representations are indexed by dominant highest weights.
In the polynomial representation ring of the general linear group,
these weights are partitions \(\lambda_1\ge\cdots\ge\lambda_r\ge0\).
Their characters are Schur polynomials \(s_\lambda(t)\).
The classical Grassmannian map produces the corresponding Schubert
classes for partitions fitting its rectangle. This symmetric-function
model is the type A construction generalized by Springer morphisms.

By contrast, put \(S_\Q=\Sym_\Q(X\otimes\Q)\). Its coordinates
\(x_i=\eps_i\) are additive linear functions on \(\mathfrak t\).
In type A they satisfy \(\sum_i x_i=0\). The basic global invariant
presentations over \(\Q\) are
\[
\begin{array}{c|l|l}
\text{type}&S_\Q^W&\text{polynomial degrees}\\ \hline
A_{m-1}&\Q[e_2(x),\ldots,e_m(x)]&2,3,\ldots,m\\
B_n,C_n&\Q[e_1(x^2),\ldots,e_n(x^2)]&2,4,\ldots,2n\\
D_n&\Q[e_1(x^2),\ldots,e_{n-1}(x^2),p]&2,4,\ldots,2n-2,n
\end{array}
\tag{5}
\]
Here \(p=x_1\cdots x_n\), and \(p^2=e_n(x^2)\).
The B/C Weyl group allows arbitrary coordinate sign changes; type D
allows only an even number. That distinction explains the Euler
invariant \(p\). Levi invariant rings combine symmetric blocks with
the relevant smaller signed-permutation group.

The three expressions \(e_i(t)\), \(e_i(x)\), and \(e_i(x^2)\) are
different. They are, respectively, a torus character, a polynomial of
degree \(i\), and a polynomial of degree \(2i\). Under the Borel map,
polynomial degree \(d\) becomes cohomological degree \(2d\).
In particular, the squared-alphabet generator \(e_i(x^2)\)
has cohomological degree \(4i\).

\section{The Borel presentation and characteristic classes}
Each weight determines an associated line bundle on \(G/B\).
Choose the line-bundle convention so that
\(\beta(\omega_i)=\eps_{s_i}\). The assignment of a weight to its first
Chern class extends multiplicatively to the Borel characteristic map.
Over \(\Q\) its full- and partial-flag presentations are
\[
\begin{split}
 H^{2*}(G/B;\Q)&\simeq S_\Q/(S_{\Q,+}^W S_\Q),\\
 H^{2*}(G/P;\Q)&\simeq
 S_\Q^{W_L}/(S_{\Q,+}^W S_\Q^{W_L}).
\end{split}
\tag{6}
\]
The subscript \(+\) means positive polynomial degree. The relations
say that positive-degree characteristic classes of global \(G\)-data
vanish on the flag. Indeed, the associated bundle of a representation
of the whole group is an ordinary trivial bundle.

For \(\Gr(r,m)\), let \(c_i\) and \(d_j\) be the Chern classes of
the dual tautological subbundle and dual quotient bundle, with the
corresponding sign convention. An integral presentation is
\[
 H^*(\Gr(r,m);\Z)=
 \frac{\Z[c_1,\ldots,c_r,d_1,\ldots,d_{m-r}]}
 {((1+c_1z+\cdots+c_rz^r)(1+d_1z+\cdots+d_{m-r}z^{m-r})=1)}.
 \tag{7}
\]
The denominator denotes all positive-degree coefficient relations.
Here \(z\) is a formal bookkeeping variable; it is not an additional
cohomology generator. The degree of \(c_i\) or \(d_i\) is \(2i\).
Eliminating one block gives the familiar complete-symmetric-function
presentation and the rectangular Schur basis.

The rational presentation (6) must not be used indiscriminately over
\(\Z\). In B/D, integral Schubert classes can require factors of two
relative to simple Chern-polynomial representatives. A rational
invariant calculation determines a rational image, not automatically
its integral lattice. Nor can one average over \(W_L\) integrally
without accounting for division by \(|W_L|\).

For actual calculations, the BGG operator
\[
 \partial_i f=\frac{f-s_i f}{\alpha_i}
 \tag{8}
\]
lowers polynomial degree by one. The braid relations make
\(\partial_w\) independent of a reduced word. For homogeneous
\(f\) of degree \(d\), its coefficient at \(\eps_w\), \(\ell(w)=d\),
is obtained by the corresponding constant divided difference, with
the chosen operator-order convention. Independently, divisor
multiplication can be computed by the Chevalley formula. Agreement
between these routes is a useful exact check; counting reduced words
as independent contributions would give the wrong answer.

\section{The Springer morphism and the published maps}
Let \(\rho:G\to\GL(V)\) have faithful differential. For simple \(G\),
the trace form
\(B_\rho(H,K)=\operatorname{Tr}(d\rho(H)d\rho(K))\)
is nondegenerate. Identify \(\mathfrak g\) with its image in
\(\operatorname{End}(V)\), and project \(\rho(g)\) orthogonally onto
this image. The resulting regular conjugation-equivariant map is
\[
 \theta_\rho:G\longrightarrow\mathfrak g.
 \tag{9}
\]
In the semisimple normalization used here,
\(\theta_\rho(1)=0\) and \(d\theta_\rho|_1=\mathrm{id}\).
Its torus restriction \(\Theta_\rho:T\to\mathfrak t\) is
\(W\)-equivariant and dominant. The construction is also the
trace-projection generalized Cayley map studied by
Kostant--Michor~\cite{KM}.

Rogers's explicit weight formula makes (9) computable~\cite{Rogers}.
Write a representation weight as the column
\(\mu=(\mu(\alpha_1^\vee),\ldots,\mu(\alpha_r^\vee))^t\), with
multiplicity \(m_\mu\), and put
\[
 M_\rho=\sum_\mu m_\mu\,\mu\mu^t,\qquad
 v_\rho(t)=\sum_\mu m_\mu e^\mu(t)\mu.
 \tag{10}
\]
If \(\Theta_\rho(t)=\sum_i a_i(t)\alpha_i^\vee\), then
\(a(t)=M_\rho^{-1}v_\rho(t)\).
The matrix entries measure the trace pairing; the vector entries
are explicit Laurent polynomials. Differentiation at the identity
provides a direct normalization check.

For the defining/vector representations of B/C/D,
\[
 \Theta_\rho^*(x_i)=\bar t_i=\frac{t_i-t_i^{-1}}2.
 \tag{11}
\]
For the defining representation of \(\SL_m\), the corresponding
formula is \(x_i\mapsto t_i-\frac1m\sum_jt_j\).
The common B/C/D formula does not identify their image lattices:
the Weyl groups, character lattices and integral Schubert geometries
are different.

Kumar defines the polynomial representation subring by
\[
 \Rep_{\rho\text{-}\poly}^{\C}(L)
   =\Theta_\rho^*\bigl(\Sym_\C(\mathfrak t^*)^{W_L}\bigr)
   \subset R(L)\otimes\C,
 \qquad
 \xi_\rho^P=\beta_\C^P\circ(\Theta_\rho^*)^{-1}.
 \tag{12}
\]
This inverse is the inverse of an injective algebra map onto its
image. It is not a claim that \(\Theta_\rho\) has a global inverse
as a morphism of varieties. Part I proves the construction and
complex surjectivity for the full Levi polynomial source
\cite[Definition 4 and Theorem 5]{Kumar}.

Part II restricts to the residual symplectic or odd orthogonal factor
in the displayed Levi, giving maps \(\xi_{n,k}\) and
\(\bar\xi_{n,k}\). Part III gives the even orthogonal map
\(\xi^D_{n,k}\)~\cite{KR,KX}. These restricted sources need not generate the whole
cohomology of the ambient flag. In the B/D notation, set
\[
 e_i=e_i(\bar h_1^2,\ldots,\bar h_k^2),\qquad
 p=\bar h_1\cdots\bar h_k\quad\text{in D},\qquad p^2=e_k.
 \tag{13}
\]
The variables \(\bar h_j=(h_j-h_j^{-1})/2\) are the tail torus
coordinates, corresponding to \(\bar t_{n-k+j}\).
The published squared-generator formula, also valid in C with its
corresponding quotient bundle, is
\[
 e_i\longmapsto
 c_i^2+2\sum_{j=1}^{i}(-1)^j c_{i+j}c_{i-j}.
 \tag{14}
\]
The \(c_j\) are the Chern classes of the indicated quotient of the
trivial defining bundle by the tautological isotropic subbundle;
\(c_0=1\), and classes outside its rank are zero.
This is Part II, Proposition 11 in C and Proposition 24 in B,
and Part III, Theorem 3.6 in D.

The stable theorems keep \(k\) fixed and let \(n\) tend to infinity.
They use a graded inverse limit of cohomology under the standard
embeddings. Part II, Theorems 16 and 29, and Part III, Theorem 3.8,
prove stable injectivity of the respective maps. Finite-rank kernels
remain: a finite-dimensional flag cannot support nonzero classes in
arbitrarily high degree. Stable injectivity and finite integral
saturation are consequently different assertions.

\section{Integral refinements: what our calculations determine}
\subsection{The source and the saturation must both be specified}
Parts II and III use useful normalized polynomial \(\Z\)-forms:
\[
 \Z[e_1,\ldots,e_k]\quad\text{in B/C},\qquad
 \Z[e_1,\ldots,e_k]\oplus p\Z[e_1,\ldots,e_k]\quad\text{in D}.
 \tag{15}
\]
The second expression is additive, with multiplication determined by
\(p^2=e_k\). These are not automatically subrings of the ordinary
integral representation ring: the halves in (11) matter.
For the stated SO/Sp tail forms, the literal character intersection
instead uses \(4^ie_i\), and in D also \(2^kp\). In polynomial
degree \(d\), its image is \(2^d\) times the normalized image.
This precise rescaling is part of the source audit~\cite{Companions}.
It must not be silently transferred to arbitrary Spin invariant
sources or arbitrary representations.

More generally, for the normalized trace construction on simply
connected simple \(G\), define
\[
 \Lambda_\rho=\{\ell\in X\otimes\Q:
                 \Theta_\rho^*\ell\in\Z[X]\},\qquad
 A_{\rho,L}=\Sym_\Z(\Lambda_\rho)^{W_L}.
 \tag{16}
\]
Our torus-source deduction identifies the full integral polynomial
source with the symmetric algebra of this linear lattice, in that
scope. Taking integral Levi invariants and taking a Borel image are
further operations: degree-one generation of the torus source does
not settle either of them. In particular, (16) records the group
isogeny and representation as part of the input.

Let \(L_d\subset H^{2d}(G/P;\Z)\) be an image of any specified integral
source and define
\[
 S_d=\Sat(L_d)=(L_d\otimes\Q)\cap H^{2d}(G/P;\Z).
 \tag{17}
\]
Integer basis changes diagonalize its inclusion with positive Smith
factors \(s_1\mid\cdots\mid s_r\). Thus
\(S_d/L_d\simeq\bigoplus_i\Z/s_i\), with unit factors retained.
The product of the factors is the index in the selected saturation.
If \(r<\rk H^{2d}\), the full ambient quotient also has a free part.
For example, the image generated by \(2h\) in projective cohomology
has factor \(2^d\) in degree \(2d\), although \(h\) and \(2h\)
generate the same rational algebra.

\subsection{A type A integral baseline}
For defining \(\SL_m\), the literal linear lattice is \(Q_{\rm rt}\).
On \(G/P_1=\mathbb P^{m-1}\), its full Levi source is generated by
the elementary symmetric functions in \(x_1-x_2,\ldots,x_1-x_m\).
Their ordinary images are \(\binom mj h^j\), and our integral
projective-family deduction gives
\[
 J_H=\bigoplus_{d=0}^{m-1}\frac{m}{\gcd(m,d)}\Z h^d,
 \qquad \gcd(m,0)=m.
\]
Thus the degree-\(2d\) factor is \(m/\gcd(m,d)\). This concerns
the literal SL source; it is not the assertion that the usual
GL polynomial-representation map fails to be surjective.

\subsection{Type C: primitivity from quaternionic geometry}
For \(a=n-k\ge1\), the normalized tail image in
\(H^*(\IG(a,2n);\Z)\) is saturated in every degree.
Its factors in degree \(4m\) are all one, indexed by partitions
of \(m\) inside the \(k\)-by-\(a\) rectangle.
This is a deduction from classical geometry, not an assertion that
the literal rescaled source is saturated.

The explanation is the compact homogeneous fibration
\[
 \frac{\Sp(n)}{U(a)\times\Sp(k)}
 \longrightarrow
 \frac{\Sp(n)}{\Sp(a)\times\Sp(k)}
       =\operatorname{HGr}(k,n).
 \tag{18}
\]
Here \(\Sp(n)\) is compact notation, whose complexification is
\(\Sp(2n)\). The fiber is \(\Sp(a)/U(a)\).
The quaternionic base has integral cohomology
\(\Z[e_1,\ldots,e_k]/(h_{a+1},\ldots,h_n)\), with \(\deg e_i=4i\);
the \(h_j\) are complete symmetric functions in the squared-root
alphabet. Its generators pull back to the squared-tail generators.
The integral Serre spectral sequence has free terms in even
bidegrees only. It collapses, and its edge image is primitive
because the other filtration quotients are free.
The ingredients are classical quaternionic Schubert theory;
see~\cite{PR,PW}. Our selected-image formulation is recorded with
its source audit in~\cite{Companions}.

The two-step type C extension uses normalized middle Chern generators
and squared-tail generators. The joint image is saturated, the sum
of the two separate images is primitive, and their intersection is
the constants. Compatible integral bases, rather than rational
dimension counts, prove these statements. They are applications of
published integral Schubert and Leray--Hirsch theory.

\subsection{Type B: a complete partition formula}
The audited type B deduction determines every factor for the normalized
map \(\bar\xi_{n,k}^{\Z}\). Put
\[
 \ell=\min(k,n-k),\quad A=\max(k,n-k),\quad
 r=\lfloor n/2\rfloor,\quad c=\lceil n/2\rceil.
\]
For \(\gamma\subseteq A^\ell\), including zero parts, define
\[
 b_i=\gamma_i+\ell-i,\qquad
 u(\gamma)=\#\{i:b_i<r\},\qquad
 E_{n,k}(\gamma)=\#\{i:b_i\ge c+u(\gamma)\}.
 \tag{19}
\]
\begin{theorem}[Repository deduction: normalized type B]
For \(1\le k<n\), the nonzero Smith factors in cohomological degree
\(4m\) are
\[
 \left\{2^{E_{n,k}(\gamma)}:
     \gamma\subseteq A^\ell,\ |\gamma|=m\right\}.
 \tag{20}
\]
The image vanishes in degrees not divisible by four.
\end{theorem}
This counts multiplicities; it does not claim that the original
Schur polynomials themselves form a Smith basis.
For instance \(n=3,k=1\) gives the factors \(1,1,2\) in degrees
\(0,4,8\), respectively. The proof constructs primitive
Schur-product prefix lattices, identifies their ranks geometrically,
and uses an adapted integral basis to recover the powers of two.
Published partition counts and character identities supply inputs;
the additional assertion is their integral-image interpretation.

\subsection{Type D: the Euler sector and integral gluing}
Define
\(\mathcal B_{n,k,e}(z)=
\sum_{\gamma\subseteq A^\ell,\ E_{n,k}(\gamma)=e}z^{2|\gamma|}\),
and use \(\mathcal B_{k,k,0}=1\), all other
\(\mathcal B_{k,k,e}=0\).
The variable \(z^d\) records cohomological degree \(2d\).
Set \(\mathcal I_{n,k}=0\) for odd \(k\). For \(k=2h\), put
\[
 \mathcal I_{n,2h}(z)=
 z^{2h(\lfloor(n-2h)/2\rfloor+1)}
 \qb{h+\lceil(n-2h)/2\rceil-1}{h}{z^2}.
 \tag{21}
\]
Here \(\qb{v}{u}{q}=\sum_{\lambda\subseteq(v-u)^u}q^{|\lambda|}\)
is the Gaussian polynomial.

\begin{theorem}[Repository deduction: normalized type D]
For \(n\ge4,\ 2\le k<n\), let \(M_D(n,k;e,d)\) count the factor
\(2^e\) in degree \(2d\). Then
\[
 \sum_dM_D(n,k;e,d)z^d
 =\mathcal B_{n,k,e}(z)+z^k\mathcal B_{n-1,k,e}(z)
   +(\delta_{e,1}-\delta_{e,0})\mathcal I_{n,k}(z).
 \tag{22}
\]
\end{theorem}
The first two terms are the square and Euler sectors. Their rational
spans are independent, but their separate integral saturations need
not sum to the saturation of the combined image. This is
\emph{integral gluing}. Every gluing direction changes
\((1,1)\) into \((1,2)\); there is no further higher-factor correction.
For \(n=4,k=2\), degree eight has factors \((1,2,2)\).
The image has rank three in an ambient rank-four group, so the
ambient cokernel also has one free direction.

These refinements answer finite-rank arithmetic questions left
undetermined by complex generator formulas and stable injectivity.
They do not replace the published theorems. The companion proofs,
dependency audits and source-normalization qualifications are
listed in~\cite{Companions}. Separately, the repository has complete
relative root-to-weight image calculations for all proper E6 and E7
parabolics. Those compare two images, not their absolute ambient
cokernels; our present priority is the classical uniform program.

\section{Equivariant cohomology: retaining the torus parameters}
The left action of \(T\) on \(G/P\) leads to the Borel construction
\[
 H_T^*(G/P;\Z)=H^*(ET\times_TG/P;\Z),\qquad
 S_\Z=H_T^*(\mathrm{pt};\Z)=\Sym_\Z X.
 \tag{23}
\]
The space \(ET\) is contractible with a free \(T\)-action, and \(BT=ET/T\).
The equivariant parameters are the first Chern classes of universal
torus characters over \(BT\). They have degree two. They are formal
coefficient classes, not a choice of a numerical torus point.
Equivariant Schubert classes form an \(S_\Z\)-basis indexed by \(W^P\).

Over \(\Q\) the presentation is
\[
 H_T^{2*}(G/P;\Q)\simeq
 S_\Q\otimes_{S_\Q^W}S_\Q^{W_L}.
 \tag{24}
\]
Write \(u\) for the left coefficient variables and \(x\) for the
right characteristic variables. The balancing relations are
\(f(x)=f(u)\) for \(f\in S_\Q^W\).
Thus the ordinary global-invariant relations \(f(x)=0\) are replaced
by parameter-dependent relations. Setting \(u=0\) recovers (6).
For integral calculations one uses an established integral
characteristic lift or a geometric presentation, rather than
assuming (24) without its rational coefficients.

\begin{example}[Projective space]
For the defining \(\SL_m\) action, choose compatible signs and write
\[
 H_T^*(\mathbb P^{m-1};\Z)
 =\Z[u_1,\ldots,u_m,h]/\left(\sum_i u_i,\prod_i(h-u_i)\right).
 \tag{25}
\]
The fixed points are the coordinate lines, and restriction sends
\(h\) to \(u_i\) at the \(i\)-th point. For \(m=2\), the relation is
\(h^2=u^2\). Forgetting equivariance gives \(h^2=0\).
\end{example}

The torus fixed points of a general \(G/P\) are indexed by \(W/W_L\).
Characteristic classes restrict by Weyl substitution; in compatible
coordinates this is \(f(w^{-1}u)\).
Restrictions along an invariant curve differ by a multiple of its
root label. After appropriate localization these restrictions
determine the class. They provide an independent route to the
same equivariant Schubert coefficients.

The equivariant extension of the original construction is
\[
 \xi_{\rho,T}^P(\Theta_\rho^*f)=\beta_T^P(f).
 \tag{26}
\]
For a specified integral source its coefficient-span image is
\(J_T^H=S_\Z\beta_T^P(A_{\rho,L})\).
Let \(I=(X)\subset S_\Z\) be the parameter augmentation ideal,
\(H_T=H_T^{2*}(G/P;\Z)\), and \(J_H\) the ordinary image.
There is a surjection \(J_T^H/IJ_T^H\to J_H\), but its injectivity is
a separate question. Since \(H_T\) is free over \(S_\Z\), its kernel is
\[
 \ker(J_T^H/IJ_T^H\longrightarrow J_H)
 \simeq \Tor_1^{S_\Z}(H_T/J_T^H,\Z).
 \tag{27}
\]
Freeness of ambient equivariant cohomology does not force freeness
of the quotient by the image.

Our integral \(\SL_4/P_1=\mathbb P^3\) example makes the issue concrete.
The defining literal source is \(\Sym_\Z(Q_{\rm rt})^{W_L}\).
Set \(a=e_2(u), b=e_3(u)\). Its elementary root-difference generators
have images
\[
 f_1=4h,\qquad f_2=6h^2+a,\qquad f_3=4h^3+2ah-b.
 \tag{28}
\]
The ordinary image is
\(\Z\oplus4\Z h\oplus2\Z h^2\oplus4\Z h^3\).
Nevertheless \(J_T^H/IJ_T^H\to J_H\) first has a \(\Z/2\) kernel
in polynomial degree five. This is an internally proved and
independently computed specialization defect~\cite{Exploration}.
It shows why a classical equivariant program must compute the
augmentation map, not only ambient ring presentations.

\section{K-theory: replacing additive classes by bundles}
The ring \(K^0(G/P)\) is generated by vector bundles, with the relation
\([E]=[E']+[E'']\) for an exact sequence and multiplication by tensor
product. A line bundle is invertible: \([L^{-1}]=[L]^{-1}\).
Unlike cohomology, ordinary \(K^0\) is not naturally graded by
codimension. It has a support, or even skeletal, filtration.
For a complex flag variety its associated graded is integral
Chow cohomology, identified with \(H^{2*}\).
This is the support filtration, not an assertion about the gamma
filtration of a representation ring.

The characteristic homomorphism is the associated-bundle map
\[
 \kappa_P:R(L)\longrightarrow K^0(G/P).
 \tag{29}
\]
For the simply connected groups in our integral setup, its
representation-ring presentation is
\[
 K^0(G/P)\simeq\Z\otimes_{R(G)}R(L).
 \tag{30}
\]
The map \(R(G)\to\Z\) is dimension. Thus a global representation
becomes an ordinary trivial bundle of its dimension. This is the
multiplicative-character counterpart of the invariant relations
in (6); it is not the same polynomial quotient.
The characteristic and Schubert frameworks are classical
\cite{GR,Karpenko}.

For projective space, take \(L=[\mathcal O(1)]\) and \(v=L-1\).
Then
\[
 K^0(\mathbb P^{m-1})=\Z[v]/(v^m),\qquad
 \ch(L)=e^h=1+h+\frac{h^2}{2!}+\cdots .
 \tag{31}
\]
The Chern character identifies K-theory with even cohomology
rationally, with higher-degree corrections. It is not generally
an integral identification in these coordinates.

The Springer-character source already supplies a canonical integral
map whenever the source characters are literal:
\[
 \eta_\rho^P:A_{\rho,L}\longrightarrow K^0(G/P),\qquad
 f\longmapsto\kappa_P(\Theta_\rho^*f).
 \tag{32}
\]
This map keeps the actual virtual character. The old map first
undoes \(\Theta_\rho^*\) and then applies \(\beta^P\).
Their relationship is a leading-filtration comparison, not the
equality of \(\ch\eta_\rho^P(f)\) with \(\beta^P(f)\) in all degrees.

On \(\Sp_4/P_1=\mathbb P^3\), for example, the normalized coordinate
\(x_1\) has integral cohomology image \(h\), but its rational
K-class is
\[
 \frac{L-L^{-1}}2=v-\frac{v^2}{2}+\frac{v^3}{2}.
 \tag{33}
\]
It is not integral. The literal generator \(2x_1\) is integral.
This is why a normalized cohomological source cannot be imported
into K-theory merely because its old Schubert coefficients were integers.

\begin{proposition}[Repository deduction: the finite filtered comparison]
For the full literal source (16), with normalized trace Springer map
on simply connected simple \(G\), the maps \(\beta^P\) and
\(\eta_\rho^P\) have the same integral source kernel.
Their images \(J_H\) and \(J_K\) are isomorphic as rings, strictly
for polynomial degree and the induced support filtration, with
\[
 \operatorname{gr}J_K=J_H,\qquad
 [K^0(G/P):J_K]=\prod_d[H^{2d}(G/P;\Z):J_H^d].
 \tag{34}
\]
\end{proposition}
\noindent
For the proof, positive global invariants pull back to virtual
global characters of dimension zero, hence vanish under \(\kappa\).
The identity first jet of \(\Theta_\rho\) makes the induced rational
map triangular with identity diagonal in the support filtration.
It is therefore invertible rationally; torsion-free integral targets
give equality of kernels on the literal source.
Homogeneous image lifts with independent leading classes establish
strictness. The determinant then gives the index identity.
The detailed argument and its qualifications are in~\cite{Exploration}.

The order of a finite cokernel does not determine its Smith factors.
Filtration layers can join by extensions, as two \(\Z/2\) layers can
form either \((\Z/2)^2\) or \(\Z/4\).
For the literal defining source on \(\Sp_4/P_1=\mathbb P^3\), put
\(b=L-L^{-1}=2v-v^2+v^3\). The image has basis \(1,b,b^2,b^3\).
Its inclusion matrix in \(1,v,v^2,v^3\) is
\[
 \begin{pmatrix}
 1&0&0&0\\0&2&0&0\\0&-1&4&0\\0&1&-4&8
 \end{pmatrix}.
\]
Its factors are \((1,1,8,8)\), whereas the ordinary cohomological
factors are \((1,2,4,8)\). Both indices are \(64\).
The type A defining-coordinate comparison does preserve the total
ordinary Smith factors for every parabolic, but that is a
representation-specific result. The completed SL3 calculations
also show that equal linear lattices for two representations need
not give equal embedded K-images~\cite{Exploration}.

\section{Equivariant K-images and the explicit quotient map}
Equivariant K-theory remembers \(T\)-actions on bundles. Its coefficient
ring is \(R(T)=\Z[X]\), rather than \(\Sym_\Z X\), and its characteristic
presentation for our simply connected setup is
\[
 K_T^0(G/P)\simeq R(T)\otimes_{R(G)}R(L).
 \tag{35}
\]
The left coefficient map sends a global representation to its
restricted torus character, not its dimension.
Forgetting equivariance is \(e^\lambda\mapsto1\).
This is distinct from setting the additive parameters of equivariant
cohomology to zero, although the augmentation-filtration comparison
connects them. Schubert structure sheaves give an \(R(T)\)-basis;
Laurent Demazure operators and fixed-point restriction formulas
provide established calculation methods~\cite{GR,Willems}.

For the Springer map, define the integral image
\[
 J_{\rho,T}^{K,P}
 =R(T)\,\kappa_T^P(\Theta_\rho^*A_{\rho,L})
       \subset K_T^0(G/P).
 \tag{36}
\]
It is this embedded submodule and subring that we seek to present.
There is already a useful rational source family.
Put \(R_\Q=\Q[X]\), \(A=S_\Q^{W_L}\), and \(B_0=S_\Q^W\). Then
\[
 \mathcal A_{\rho,T}^P
 =R_\Q\otimes_{B_0,\,f\mapsto f(\Theta_\rho(t))}A
 \longrightarrow K_T^0(G/P)_\Q,\qquad
 1\otimes f\longmapsto\kappa_T^P(\Theta_\rho^*f).
 \tag{37}
\]
The map from \(B_0\) is essential: the basic invariant values
\(f_i(\Theta_\rho(t))\) are exactly the coordinates of the induced
adjoint-quotient map \(T/W\to\mathfrak t/W\).
The previous quotient-coordinate computations therefore become
coefficients of relations in a characteristic-image problem.
At \(t=1\) these positive-degree coefficients vanish, giving the
rational ordinary source quotient. A general integral freeness
statement is not being inferred from (37).

\subsection{A determinant theorem and its geometric meaning}
For each root-length orbit \(a\), define
\[
 \Delta_a(x)=\prod_{\alpha>0,\ {\rm length}=a}\alpha(x),\quad
 \varpi_a=\tfrac12\sum_{\alpha>0,\ {\rm length}=a}\alpha,\quad
 \delta_a(t)=e^{\varpi_a}\prod_{\alpha>0,\ {\rm length}=a}(1-e^{-\alpha}).
\]
The root quotient
\(\mathcal J_{\rho,a}=\Delta_a(\Theta_\rho(t))/\delta_a(t)\)
is a regular invariant Laurent polynomial over \(\Q\), with value
one at the identity. Let \(N=|W/W_L|\), and let \(N_a,N_{L,a}\)
count positive roots of that length in \(G,L\). Define
\[
 m_{P,a}=\frac N2\left(1-\frac{N_{L,a}}{N_a}\right).
 \tag{38}
\]
This integer is half the number of cosets moved by a reflection
of that length.

\begin{theorem}[Repository deduction over \(R_\Q\)]
The map (37) is injective between finite free \(R_\Q\)-modules of
rank \(N\). If \(C_{\rho,T}^P\) is its cokernel, then
\[
 \Fitt_0(C_{\rho,T}^P)
   =\left(\prod_a\mathcal J_{\rho,a}^{\,m_{P,a}}\right).
 \tag{39}
\]
\end{theorem}
The proof compares two fixed-point matrices. A homogeneous polynomial
basis evaluates as \(f_j(w^{-1}\Theta_\rho(t))\); each reflection
wall forces coincident row pairs. Divisibility and the total degree
determine its root-product determinant. The target Schubert
restriction matrix has the torus Euler root products as determinant.
Their ratio is (39), up to a unit.

For a presentation \(R^b\to R^a\to C\to0\),
\(\Fitt_i(C)\) is generated by the \((a-i)\)-minors.
Thus \(\Fitt_0\) is the determinant ideal in the square case;
higher Fitting ideals record finer degeneracy information.
The rings \(\Z[X]\) and \(\Sym_\Z X\) are multivariable rings, so
there is no general Smith normal form over them.
Smith factors are appropriate after passing to a specified DVR
or to a finite integer-lattice specialization.

Write \(\mathcal J_\rho=\prod_a\mathcal J_{\rho,a}\).
The full Springer differential determinant has torus restriction
\[
 D_\rho|_T=\det(d\Theta_\rho)\,\mathcal J_\rho^2.
 \tag{40}
\]
Consequently (39) detects the root-direction factors of the generalized
Cayley determinant. The Cartan Jacobian is a separate factor; these
image calculations do not by themselves determine the full critical
divisor, target branch locus, or generic degree.

\subsection{Two arbitrary-rank classical families already available}
For the defining \(\SL_m\)-module \(V\), the rational equivariant map
(37) is an isomorphism for every \(P\).
The trace formula gives
\(\Theta_{V^{\otimes j}}=(\chi_V/m)^{j-1}\Theta_V\).
A homogeneous basis of degrees \(\ell(w)\) therefore yields
the complete rational cokernel decomposition
\[
 C_{V^{\otimes j},T}^P
 \simeq\bigoplus_{w\in W^P}
       R_\Q/(\chi_V^{(j-1)\ell(w)}),\qquad j\ge1.
 \tag{41}
\]
A term with exponent zero is the zero module.
This supplies every higher Fitting ideal in this family, not just
the determinant. The literal integral source remains an additional
arithmetic input.

There is also a complete integral source model for defining
\(\Sp_{2n}/P_1=\mathbb P^{2n-1}\).
Put
\[
 R=\Z[t_1^{\pm1},\ldots,t_n^{\pm1}],\quad
 s_i=t_i+t_i^{-1},\quad B_i=t_i-t_i^{-1},\quad b=L-L^{-1}.
\]
Then
\[
 K_T^0(\mathbb P^{2n-1})
 =R[L]/\prod_i(L-t_i)(L-t_i^{-1}),\qquad
 J_{\rho,T}^{K,P_1}=R[b]
 \simeq R[z]/\prod_i(z^2-B_i^2).
 \tag{42}
\]
The omitted tail source generators are eliminated by the global
squared-symmetric relations. The image is free on
\(1,b,\ldots,b^{2n-1}\), and the integral determinant is
\[
 \Fitt_0\bigl(K_T^0/J_{\rho,T}^{K,P_1}\bigr)
   =\left(2^n\prod_{i<j}(s_i+s_j)^2\right).
 \tag{43}
\]
To see the factors, compare the Vandermonde matrices for
\(L\) and \(L-L^{-1}\) at
\(t_1,t_1^{-1},\ldots,t_n,t_n^{-1}\).
Each reciprocal pair contributes two; the four cross-pairs
for \(i<j\) contribute \((s_i+s_j)^2\).
At \(t_i=1\), the determinant is \(2^{n(2n-1)}\).
This family's image specializes faithfully to its ordinary K-image.
Its arbitrary-rank higher Fitting ideals remain a substantive next
problem. Proofs and independent finite checks of (34), (39),
(41)--(43) are recorded in~\cite{Exploration}.

\section{A classical program with precise next goals}
The common goal is to determine the embedded integral image (36),
its module presentation, its Fitting ideals and its augmentation
to ordinary K-theory, for the defining/vector representation in
the classical families. Equivariant ordinary cohomology is a parallel
comparison target. The determinant theorem is an established starting
point within the companion notes, not the remaining classification.

Type A supplies a controlled baseline: construct the actual integral
source image on Grassmannians and general partial flags, retaining
the difference between the root and weight lattices. The rational
answer alone cannot see the integral specialization defects of
Section~7. Type C next asks for all higher Fitting ideals in (42),
and then for isotropic Grassmannians with subspace dimension greater
than one. The known quaternionic arguments explain ordinary
primitivity, but do not automatically give this equivariant
multiplicative classification.

For B and D, begin with isotropic lines. The targets are the odd
quadric \(\Spin_{2n+1}/P_1\) of dimension \(2n-1\) and the even
quadric \(\Spin_{2n}/P_1\) of dimension \(2n-2\).
Their K-theory must retain the full Schubert or spinor-bundle data.
In D the two middle Schubert classes must both be retained.
The vector representations factor through SO and have faithful
differential on Spin; the coefficient lattice still records which
group is chosen. A polynomial projective-bundle model cannot simply
be substituted for a quadric target.

These first models should lead to isotropic Grassmannians and then
arbitrary partial flags. Type D maximal isotropic components must
be specified separately. Comparisons involving adjoint, tensor,
exterior-power and spin representations are later tests of which
features depend only on the linear source and which depend on the
full Springer character.

An effective computational implementation should return the source
lattice and normalization, the coefficient ring, the actual image
matrix or relations, Fitting data, and the augmentation map.
Verification should compare fixed-point restrictions with
Schubert/Demazure calculations, and recover the appropriate old
ordinary image with its stated source. Exact small-rank checks test
a proposed uniform argument; they do not replace its proof.

The resulting picture extends the maps of Parts I--III in a precise
way. Their polynomial characters and Borel images remain the source
of the construction. The integral work asks which lattice of
cohomology those characters generate. Equivariance retains the torus
parameters that ordinary cohomology discards; K-theory retains the
multiplicative character and extension information that the leading
Chern-class calculation discards. The unfinished task is to explain
these integral images uniformly across the classical root data.

\begin{thebibliography}{99}
\small\raggedright\setlength{\itemsep}{2pt}
\bibitem{Kumar}
S.~Kumar, \emph{Representation ring of Levi subgroups versus cohomology
ring of flag varieties}, Math.\ Ann.\ \textbf{366} (2016), 395--415.
\href{https://kumar.math.unc.edu/papers/kumar76.pdf}{Author's text}.
\bibitem{KR}
S.~Kumar and S.~Rogers, \emph{Representation ring of Levi subgroups versus
cohomology ring of flag varieties II}, J.\ Algebra \textbf{556} (2020),
340--362. \href{https://kumar.math.unc.edu/papers/kumar81.pdf}{Author's text}.
\bibitem{KX}
S.~Kumar and J.~Xie, \emph{Representation ring of Levi subgroups versus
cohomology ring of flag varieties III}, J.\ Pure Appl.\ Algebra
\textbf{228} (2024), 107539.
\href{https://kumar.math.unc.edu/papers/kumar86.pdf}{Author's text}.
\bibitem{Rogers}
S.~Rogers, \emph{An explicit determination of the Springer morphism},
Comm.\ Algebra \textbf{46} (2018), 4233--4242.
\href{https://arxiv.org/abs/1701.01538}{Available preprint}.
The preprint is distinct from the final journal text.
\bibitem{Brion}
M.~Brion, \emph{Lectures on the geometry of flag varieties},
\href{https://arxiv.org/abs/math/0410240}{arXiv:math/0410240}.
\bibitem{PR}
P.~Pragacz and J.~Ratajski, \emph{Formulas for Lagrangian and orthogonal
degeneracy loci; \(\widetilde Q\)-polynomial approach},
Compositio Math.\ \textbf{107} (1997), 11--87.
\href{https://arxiv.org/abs/alg-geom/9602019}{Preprint}, especially Appendix A.
\bibitem{PW}
I.~Panin and C.~Walter, \emph{Quaternionic Grassmannians and Borel classes
in algebraic geometry},
\href{https://arxiv.org/abs/1011.0649}{arXiv:1011.0649}, Theorem 11.2.
\bibitem{GR}
S.~Griffeth and A.~Ram, \emph{Affine Hecke algebras and the Schubert
calculus}, \href{https://arxiv.org/abs/math/0405333}{arXiv:math/0405333},
especially Sections 1--4.
\bibitem{Willems}
M.~Willems, \emph{K-th\'eorie \'equivariante des vari\'et\'es de drapeaux
et des vari\'et\'es de Bott--Samelson},
\href{https://arxiv.org/abs/math/0204265}{arXiv:math/0204265}.
\bibitem{Karpenko}
N.~Karpenko, \emph{A remark on connective K-theory},
\href{https://sites.ualberta.ca/~karpenko/publ/connective.pdf}{Author's text},
especially Theorems 2.2 and 3.1.
\bibitem{KM}
B.~Kostant and P.~W.~Michor, \emph{The generalized Cayley map from an
algebraic group to its Lie algebra}, in \emph{The Orbit Method in Geometry
and Physics}, Progress in Mathematics \textbf{213} (2003), 259--296.
\href{https://arxiv.org/abs/math/0109066}{arXiv:math/0109066}.
\bibitem{Companions}
\emph{Integral Springer--Borel companion proofs}, unpublished repository
notes, internally audited October 2026:
\record{TRACE_SOURCE_THEOREM.md}{torus source theorem};
\record{INTEGRAL_LITERATURE.md}{normalization and source audit};
\record{TYPE_B_GENERAL_SMITH.md}{type B multiplicities};
\record{TYPE_D_EVEN_GLUING_THEOREM.md}{type D gluing};
\record{TYPE_C_JOINT_AUDIT.md}{type C proof and literature audit};
\record{PROOF_STATUS.md}{proof status}.
These deductions have not received external expert review, and priority
for their exact formulations has not been established.
\bibitem{Exploration}
\emph{Other cohomology theories and affine groups}, unpublished repository
exploration, 22 October 2026:
\record{OTHER_THEORIES_FINITE.md}{finite K comparison};
\record{OTHER_THEORIES_SL3.md}{SL3 formulas};
\record{OTHER_THEORIES_TYPE_A.md}{type A coordinate comparison};
\record{EQUIVARIANT_IMAGE_PILOT.md}{integral specialization defect};
\record{OTHER_THEORIES_EQUIVARIANT_K.md}{equivariant K source and determinants};
\record{other-theories-verification.json}{bound calculation reports}.
The general arguments and finite independent checks have different scopes;
neither is a certificate of originality.
\end{thebibliography}
\end{document}
